\section{Поисковая репликация алгоритма INoT}
\label{sec:inot-replication}
Четыре факторных условия проверяют отдельные компоненты, но не воспроизводят опубликованный метод целиком. Поэтому отдельно зафиксирована промптовая репликация INoT~\cite{sun2025inot}, обозначенная INoT*. Она использует те же 200 задач, полный предоставленный контекст, GPT-5.4 mini medium, итоговый формат и контроль оцениваемости. На задачу приходится один новый вызов, работают два работника с тем же пределом 180 секунд. Независимо сформулированная инструкция PromptCode описывает двух виртуальных участников, исходные ответы, взаимные аргументы, критику, возражения, обновления и проверку согласия с пределом десять раундов. Это концептуальная инструкция для модели, а не код, исполняемый на компьютере, или проверенная последовательность скрытых взаимодействий агентов.

Правило остановки следует текстовому описанию исходной статьи: согласие либо предел раундов. Соответствующий исходный листинг содержит неоднозначное условие цикла и не задаёт итоговый ответ при отсутствии согласия. В этой репликации явно возвращается последний ответ виртуального участника A. Дополнение изображениями исключено для текстовых задач. Точная инструкция приведена в приложении с запросами. Проверенная исполняемая авторская реализация для данной постановки не найдена; подтверждение адаптации авторами не заявляется.

\begin{figure}[htbp]\centering
\begin{tikzpicture}[x=1mm,y=1mm,
 box/.style={draw=black!80,align=center,font=\small,inner sep=2mm},
 flow/.style={-{Latex[length=2mm]},thick}]
\node[box,fill=gray!8,text width=99mm] (input) at (0,0) {Один запрос к модели: задача, контекст и инструкция INoT*};
\draw[dashed,black!65] (-52,-12) rectangle (52,-114);
\node[font=\small,align=center,text width=100mm] at (0,-18) {\textbf{Содержание инструкции, не наблюдаемая трасса}};
\node[box,text width=77mm,minimum height=15mm] (initial) at (0,-36) {Предложить решения A и B\\и их обоснования};
\node[box,text width=77mm,minimum height=19mm] (exchange) at (0,-62) {Обменяться аргументами и критикой;\\ответить на критику; обновить решения};
\node[diamond,draw=black!80,aspect=2.7,align=center,font=\small,inner sep=1mm,text width=34mm] (decision) at (0,-95) {Ответы совпали\\или завершены\\10 раундов?};
\node[box,fill=gray!8,text width=99mm,minimum height=12mm] (answer) at (0,-129) {Вернуть последнее решение A};
\draw[flow] (input.south) -- (0,-12);
\draw[flow] (0,-24) -- (initial.north);
\draw[flow] (initial) -- (exchange);
\draw[flow] (exchange) -- (decision);
\draw[flow] (decision.west) -- (-47,-95) node[above right,font=\small] {нет} |- (exchange.west);
\draw[flow] (decision.south) -- node[right,pos=0.2,font=\small,fill=white,inner sep=1pt] {да} (answer.north);
\end{tikzpicture}

\caption{Как читать инструкцию INoT*. Пунктирная область описывает действия, предписанные в одном запросе: это не исполняемый цикл и не измеренные внутренние агенты. При согласии ответ A совпадает с ответом B; без согласия после десяти раундов инструкция выбирает последний ответ A. Дословный текст приведён в приложении~\ref{app:inot-instruction}.}
\label{fig:inot-instruction}
\end{figure}
\FloatBarrier

Независимое административное продолжение использует то же правило «только ещё не отправленные»: начальная попытка с превышением времени не повторяется, а оставшиеся 137 нетронутых назначений фиксируются и отправляются по одному разу. Все трассы первой серии и продолжения хранятся раздельно. Сравнения INoT* с D и SR являются поисковыми и не входят в четыре основных теста. Парность контролирует состав задач, но отдельная серия сохраняет различия времени отправки, состояния кеша и возможного изменения модели у поставщика. Наблюдаемый результат характеризует эту раскрытую адаптацию, а не точное воспроизведение чисел исходной статьи.
